\chapter{Research program and paper architecture}\label{ch:program}
\index{open problem}
\index{research program}


\section{The central research problem}
For a nested pair of gate models, determine the distortion function
\[
D_n=
\sup_{U\in\mathcal T_n}
\bigl(c_{\mathrm{res}}(U)-c_{\mathrm{amb}}(U)\bigr),
\]
where $\mathcal T_n$ is the target class.  The basic questions are:
\begin{enumerate}
\item Does any defect exist?
\item What is the first size threshold?
\item What is the minimal catalytic motif?
\item Is the effect generic, sparse, or confined to special strata?
\item How does it scale under direct sum, recursion, and tensor products?
\item Which resource vector reveals the first defect?
\item Can the effect survive realistic architecture and fault-tolerant costs?
\end{enumerate}

\section{Established base camp}
The unitriangular theory provides a complete first example:
\begin{itemize}
\item a universal six-row atom over nonzero commutative rings;
\item exact binary and ternary first-threshold evidence;
\item exact classification of twenty binary minimal defects;
\item connected families with two reverse gates and linear saving;
\item microscopic rank-one, bandwidth-one escape;
\item generic rigidity and defect codimension;
\item recursive and tensor amplification to $\Omega(n\log n)$.
\end{itemize}
The small Clifford atlas adds a second gate logic with genuine exact defects and shows that changing the primitive library by adding SWAP can remove them.

\section{Program A: complete the Clifford threshold theory}
The highest-priority quantum project is a paper on parabolic-versus-ambient Clifford synthesis.

\subsection{Immediate tasks}
\begin{enumerate}
\item Classify the four two-qubit entangling defects under local Clifford and qubit-swap equivalence.
\item Classify the twenty-four three-qubit total-count defects under parabolic conjugation and qubit permutations.
\item Derive analytic normal forms for the witness circuits.
\item Prove lower bounds independent of full group enumeration.
\item Test depth, weighted, and architecture-aware metrics.
\item Determine whether a finite defective atom tensor-amplifies.
\end{enumerate}

\subsection{Longer-term questions}
\begin{conjecture}
The parabolic subgroup of $\Sp(2n,2)$ has unbounded additive distortion in the full symplectic group for at least one natural Clifford cost model.
\end{conjecture}
\begin{question}
Does the maximum total-count defect grow as $n$, $n\log n$, or faster?
\end{question}
\begin{question}
Is the disappearance of three-qubit defects after adding unit-cost SWAP a small-width accident or a structural geodesic-convexity phenomenon?
\end{question}

\section{Program B: temporary nonlinearity}
The three-bit affine-versus-Toffoli census is rigid.  The next target is a complete four-bit normal-form search without enumerating $16!$ permutations.

\subsection{Algorithmic plan}
\begin{enumerate}
\item Compute exact affine distances for all $|\operatorname{AGL}(4,2)|=322{,}560$ targets.
\item Enumerate Toffoli sandwiches modulo affine left and right actions.
\item Use ANF equations to retain only affine endpoints.
\item Compare total, Toffoli, CNOT, and fault-tolerant costs.
\item Produce SAT certificates below every claimed optimum.
\end{enumerate}

\begin{question}
What is the first bit width at which a restored nonlinear excursion shortens an affine target?
\end{question}
\begin{question}
Can temporary degree two yield a family with unbounded saving while the endpoint degree remains one?
\end{question}

\section{Program C: phase-polynomial catalysts}
The phase setting is likely to produce a defect first in a resource other than total primitive count.

Priority metrics:
\[
T\text{-depth},
\quad
\text{CNOT count at fixed }T\text{-count},
\quad
\text{two-qubit depth},
\quad
\text{architecture-weighted cost}.
\]
A first paper should choose one metric and avoid a diffuse comparison across all resources.

\begin{question}
Can a restored Hadamard or Clifford frame reduce $T$-depth of a diagonal target without additional ancillas?
\end{question}
\begin{question}
Does tensor activation occur for phase-polynomial $T$-count?
\end{question}

\section{Program D: Gaussian and matchgate systems}
Two concrete first searches are recommended.
\begin{enumerate}
\item Number-conserving versus general fermionic Gaussian synthesis on a small number of modes.
\item Nearest-neighbor versus all-to-all matchgate circuits with weighted fermionic SWAP cost.
\end{enumerate}
Both admit exact orthogonal-matrix representations and have clear architecture relevance.

\section{Program E: symmetry and measurement}
This program should begin with a narrowly declared resource theory, such as:
\begin{itemize}
\item excitation-preserving two-qubit gates versus arbitrary two-qubit gates;
\item unitary Clifford+$T$ versus measurement-assisted Clifford+$T$;
\item clean versus dirty ancillas for a fixed arithmetic family;
\item local architecture versus one reusable nonlocal bus gate.
\end{itemize}
The first objective is not maximal generality but one theorem with exact restoration and cost accounting.

\section{Program F: root-system catalysis}
For a split reductive group with allowed roots $\Psi\subset\Phi$, classify configurations in which a forbidden root acts as a reusable catalyst through Chevalley commutator relations.

Suggested sequence:
\[
A_2\to B_2=C_2\to G_2\to A_r.
\]
Low-rank nonsimply laced systems are attractive because longer root strings may create more than one useful effect per consumer.

\begin{question}
Is there a root-system analogue of the three-cross, one-phase threshold?
\end{question}
\begin{question}
Can a $G_2$ root string produce a multiplicative distortion exceeding the clean type-$A$ value $3/2$?
\end{question}

\section{Program G: algebraic defect loci}
For each finite-dimensional gate hierarchy, define bounded word loci and defect strata.  Determine:
\begin{itemize}
\item irreducible components,
\item exact dimension and codimension,
\item singularities and degenerations,
\item border length,
\item finite-field point counts,
\item behavior under tensor products.
\end{itemize}

\begin{conjecture}
In a broad class of one-parameter elementary gate varieties, defect at least $d$ is generically a codimension-at-least-$d$ condition, with equality on explicit families of components.
\end{conjecture}

\section{Program H: packet-resolved realization}
The packet hierarchy should be developed as a separate mathematical framework rather than appended to every circuit paper.

Core objectives:
\begin{enumerate}
\item define spherical, type, and holonomy functors for finite transition groupoids;
\item prove information-loss theorems for commutativization;
\item classify when typed Green kernels are sufficient for distance;
\item relate imprimitivity to combinatorial decomposition;
\item formalize realization fibers and descent obstructions.
\end{enumerate}

\section{Recommended paper sequence}
\begin{longtable}{@{}>{\raggedright\arraybackslash}p{.08\textwidth}>{\raggedright\arraybackslash}p{.37\textwidth}>{\raggedright\arraybackslash}p{.45\textwidth}@{}}
\toprule
Paper & Working title & Core result\\
\midrule
I & Catalytic escape from unitriangularity & Completed linear threshold, rarity, and asymptotic distortion.\\
II & Catalytic basis escape in Clifford synthesis & Exact parabolic-versus-Clifford thresholds and analytic motifs.\\
III & Temporary nonlinearity in reversible circuits & Affine targets with or without Toffoli-assisted defects.\\
IV & Gate-set distortion and catalytic holonomy & Abstract subgroup, groupoid, and vector-cost theory.\\
V & Spherical collapse and transition holonomy & Packet refinement, Latin-square holonomy, and realization fibers.\\
VI & Algebraic defect loci & Word varieties, codimension, border catalysis, and point counts.\\
VII & Root-system catalytic complexity & Chevalley commutators and thresholds beyond type $A$.\\
\bottomrule
\end{longtable}

\section{Software roadmap}
Version $0.1$ of the atlas supports exact small finite groups.  A mature package should add:
\begin{enumerate}
\item generic multiobjective Dijkstra search;
\item disk-backed and parallel BFS;
\item SAT and SMT export;
\item group action and orbit classification;
\item symbolic word-map elimination;
\item circuit diagrams in Quantikz and SVG;
\item architecture graphs and routing costs;
\item certificate verification independent of the search engine;
\item a machine-readable theorem-status ledger.
\end{enumerate}

\section{Standards for claims}
A result should be labeled \emph{proved} only when every lower bound and restoration claim has a proof.  A complete finite census should identify the state quotient, enumerate its full size, and make the traversal reproducible.  A symbolic skeleton should be valued as a research object but not described as an established shortcut.

\section{Potential applications}
If developed, the theory may contribute to:
\begin{itemize}
\item exact and heuristic quantum circuit compilers,
\item architecture-aware routing,
\item parity and phase-network synthesis,
\item stabilizer and Clifford optimization,
\item reversible arithmetic,
\item Gaussian and fermionic simulation,
\item lower-bound methodology,
\item geometric group theory of gate sets,
\item algebraic geometry of circuit word maps.
\end{itemize}
The near-term value is methodological and algorithmic.  Hardware benefit requires a device-specific cost and noise model.

\section{Twenty central open problems}
\begin{enumerate}
\item Prove or refute a universal six-dimensional linear threshold over all coefficient rings.
\item Determine the true asymptotic order of the maximum unitriangular defect.
\item Decide whether multiplicative distortion is bounded.
\item Determine the lower border length of the canonical six-row atom.
\item Classify tree-tight catalytic matroids.
\item Find a one-copy rigid target with tensor activation.
\item Analytically classify the two-qubit Clifford entangling defects.
\item Analytically classify the three-qubit total-count Clifford defects.
\item Determine whether SWAP restores rigidity for all $n$ in the chosen primitive model.
\item Find the first affine target shortened by temporary nonlinearity, or prove none exists.
\item Find the first diagonal target improved by non-diagonal Clifford catalysis.
\item Determine a phase-polynomial Pareto defect.
\item Find a number-conserving Gaussian target shortened by Bogoliubov escape.
\item Find a matchgate routing catalyst with a certified architecture advantage.
\item Establish a symmetry-sector microscopic-escape theorem.
\item Classify when a fixed number of dirty ancillas gives unbounded saving.
\item Develop typed Green-kernel lower bounds that avoid full enumeration.
\item Determine finite moment bounds sufficient for transition-holonomy classification.
\item Classify catalytic motifs in $B_2$, $C_2$, and $G_2$.
\item Prove a general theorem connecting algebraic codimension and catalytic defect.
\end{enumerate}

\section{Closing perspective}
The common phenomenon is not that every restriction can be beaten.  It is that the path space matters.  A target may obey a symmetry, flag, degree bound, locality rule, or factorization constraint while the cheapest route to that target does not.

The disciplined research question is therefore:
\begin{center}
\fbox{\parbox{.90\textwidth}{\centering
Which temporary resources can be opened, reused, and exactly removed at a net computational gain?}}
\end{center}
The established linear theory answers this question completely for one important model.  The monograph supplies the symbolic and computational machinery to ask it across quantum gate logic without conflating identities, experiments, and theorems.
